% !PRINTED 64
\weberpage{64}

\begin{center}
{\large\gesperrt{Second Section.}}\\[0.5em]
{\Large\gesperrt{Determinants.}}\\[0.4em]
\end{center}

\begin{center}
{\large \S.\ 18.}

\vspace{0.5em}
{\large\gesperrt{Permutations of \(n\) Elements.}}
\end{center}

We consider a system of \(n\) distinct elements of any kind, for example the \(n\) digits
\[
 1,2,3,\ldots,n,
\]
whose complex, in this definite arrangement, we shall denote by \(\mathfrak A\). The elements of \(\mathfrak A\) can be arranged in different ways, for example
\[
 2,1,3,\ldots,n.
\]
The passage from one arrangement to another is called a \gesperrt{permutation}.

If we denote by \(\Pi(n)\) the number of the different arrangements, which can depend only on the number \(n\) of the elements, then first
\[
 \Pi(1)=1,\qquad \Pi(2)=2.
\]
To determine the number in general, we think of an \(n\)-th element as added to \(n-1\) elements. In every arrangement of the \(n-1\) elements the \(n\)-th element can now be put in \(n\) different places, namely before the first, between the first and the second, between the second and the third, and so on, finally after the \((n-1)\)-st; and all arrangements thus obtained are distinct. Hence
\begin{equation}
 \Pi(n)=n\Pi(n-1),
\tag{1}
\end{equation}
from which, by complete induction,
\begin{equation}
 \Pi(n)=1\cdot2\cdot3\cdots n
\tag{2}
\end{equation}
%CONT>>
